#MODELS EXPLAINATION
Yeh Basically Ek Employee Management System Hai Jisme
▪️ User ➡️ Base Entity Hai
▪️ EmployeeProfile ➡️ Extra Details
▪️ Department ➡️ Grouping
▪️ Task / Leave / Attendace ➡️ Activity Tracking

🔥 1. One-to-One Relationship (1:1)
📍 Code: user = models.OneToOneField(User, on_delete=models.CASCADE)

💡 Meaning:
➡️  Ek User Sirf Ek EmployeeProfile
➡️  Ek EmployeeProfile Ka Sirf Ek User

🧠 Real Understanding:
Har employee ke paas ek account (User)
Aur us account ka ek hi profile

👉 Duplicate profile allowed nahi ❌
👉 “User and EmployeeProfile have One-to-One relationship because each user has exactly one profile.”

🔥 2. One-to-Many Relationship (1:N)

✅ (A) Department → EmployeeProfile
📍 Code: department = models.ForeignKey(Department, ...)

💡 Meaning:

👉 Ek Department → multiple employees
👉 Ek Employee → ek hi department

🧠 Real Example:
IT Department → 10 employees
HR Department → 5 employees

👉 “One department can have many employees, but each employee belongs to only one department.”

✅ (B) User → Task
📍 Code:
        assigned_to = models.ForeignKey(User, related_name="tasks")
        assigned_by = models.ForeignKey(User, related_name="created_tasks")


💡 Meaning:

👉 Ek user → multiple tasks assign ho sakte
👉 Ek user → multiple tasks create kar sakta

🧠 Real Example:
Manager → 20 tasks assign kare
Employee → 5 tasks assigned

| Field       | Meaning          |
| ----------- | ---------------- |
| assigned_to | Task kis ko diya |
| assigned_by | Task kisne diya  |

👉 Isko bolte hai self-referencing ForeignKey

👉 “User has a one-to-many relationship with Task for both assignment and creation 
    using self-referencing foreign keys.”

✅ (C) User → Leave
employee = models.ForeignKey(User)

👉 Ek user → multiple leave apply kar sakta

✅ (D) User → Attendance
employee = models.ForeignKey(User)

👉 Ek user → daily attendance (multiple records)

🔥 3. Many-to-Many Relationship (M:N)

❌ Tere current project mein direct ManyToManyField nahi hai

💡 BUT Important 🔥 (Interview Trick)

Agar tu bole:

👉 “User ↔ Task many-to-many hai?”

Answer:
❌ Nahi

👉 Kyunki:

Ek task → sirf ek assigned_to
Isliye yeh One-to-Many hai
💡 Agar future mein banana ho:

Example:

assigned_to = models.ManyToManyField(User)

👉 Tab:

Ek task → multiple users
Ek user → multiple tasks

🧠 unique_together in Attendance
unique_together = ['employee', 'date']

💡 Meaning:

👉 Ek user ka ek din mein sirf 1 attendance record

👉 “This ensures that an employee cannot have duplicate attendance records for the same date.”


💣 Full Ready OTP And Email Answer (Practice This 🔥)


“Sir, maine OTP-based password reset system banaya hai jisme user email enter karta hai aur system OTP generate karke email par send karta hai.
Pehle system email verify karta hai, fir 6-digit OTP generate karke database mein timestamp aur attempt count ke saath store karta hai.
Maine 60 seconds ka cooldown bhi add kiya hai taaki multiple requests na aaye.
OTP ko reusable email function ke through HTML template ke saath send kiya jata hai.
Verification ke time OTP expiry (5 minutes), attempt limit aur OTP match check hota hai.
Agar valid ho to password securely update hota hai.
Overall maine system ko secure banane ke liye expiry, cooldown aur attempt limiting implement ki hai.”